\documentclass{article}
\usepackage[top=2cm, bottom=2cm, left=1.5cm, right=1.5cm]{geometry}
\usepackage{import}
\subimport{header/}{header}

\title{Magnetic pressure and tension}
\author{Your Name}
\date{}

\begin{document}
\maketitle

\section{Decomposing the Lorentz force}
\label{sec:magnetic-tension}

In MHD, the Lorentz force per unit volume\footnote{
    We adopt rationalised-Gaussian units, such that the current density is $\mVector{j} = \mCurl{\mVector{b}}$ (\ie no factor of $4\pi$).
} is
\begin{align}
    \mVector{f}_L = \rbrac{\mCurl{\mVector{b}}}\times\mVector{b}
    , \label{eqn:magnetic-tension:lorentz}
\end{align}
where $\mVector{b}$ is the magnetic field. Using the identity $\rbrac{\mCurl{\mVector{b}}}\times\mVector{b} = \rbrac{\mVector{b}\cdot\nabla}\mVector{b} - \mGrad{\rbrac{b^2/2}}$,
\begin{align*}
    \mVector{f}_L
    = \rbrac{\mVector{b}\cdot\nabla}\mVector{b}
        - \mGrad{\rbrac{\frac{b^2}{2}}} ,
\end{align*}
\ie \cref{eqn:magnetic-tension:lorentz} is the divergence of the Maxwell stress tensor,
\begin{align}
    \mVector{f}_L = \mDiv{\mTensor{T}} ,
    \qquad
    \mTensor{T} \equiv \mVector{b}\otimes\mVector{b} - \frac{b^2}{2} \imatrix
    . \label{eqn:magnetic-tension:stress}
\end{align}

\subsection{Pressure and tension}
\label{sec:magnetic-tension:decomposition}

To see the physical processes hidden within \S\ref{sec:magnetic-tension}, we introduce $\mVectorUnit{b} = \mVector{b} / b$ and decompose $\rbrac{\mVector{b}\cdot\nabla}\mVector{b}$ along both the field itself, $\mVectorUnit{b}$, and its curvature vector $\mVector{\kappa} \equiv \rbrac{\mVectorUnit{b}\cdot\nabla}\mVectorUnit{b}$ (which points in the direction the field line bends), with $\mVectorUnit{\kappa} = \mVector{\kappa} / \kappa$ and $\kappa \equiv \nbrac{\mVector{\kappa}}$. This gives,
\begin{align}
    \rbrac{\mVector{b}\cdot\nabla}\mVector{b}
    = b \mppFrac{b}{s} \mVectorUnit{b}
        + b^2\kappa \mVectorUnit{\kappa}
    , \label{eqn:magnetic-tension:decomposition}
\end{align}
where $s$ is arc length along the field line. Substituting \cref{eqn:magnetic-tension:decomposition} into \cref{eqn:magnetic-tension:lorentz} gives
\begin{align}
    \mVector{f}_L
    &= -\mGrad{\rbrac{\frac{b^2}{2}}}
        + b^2\kappa \mVectorUnit{\kappa}
        + b \mppFrac{b}{s} \mVectorUnit{b}
    \nonumber \\
    &= b^2 \mVector{\kappa}
        + \frac{1}{2} \rbrac{\mVectorUnit{b}\otimes\mVectorUnit{b}\cdot\mGrad{\rbrac{b^2}}}
        - \frac{1}{2} \mGrad{\rbrac{b^2}}
    \nonumber \\
    &= b^2 \mVector{\kappa}
        - \wrapunderbrace[3cm]{\frac{1}{2} \mGradPerp{\rbrac{b^2}}}{pressure gradient orthogonal to the field}
    . \label{eqn:magnetic-tension:final}
\end{align}
with $\mGradPerp{\rbrac{b^2}} \equiv \rbrac{\imatrix - \mVectorUnit{b}\otimes\mVectorUnit{b}}\cdot\mGrad{\rbrac{b^2}}$. The first term of \cref{eqn:magnetic-tension:final} is the \emph{magnetic tension}: it points along $\mVectorUnit{\kappa}$, is proportional to the field-line curvature $\kappa$, and vanishes for straight field lines, acting like a taut string pulling bent field lines straight. The second is the \emph{magnetic pressure} gradient, but only the component perpendicular to the field; the component along $\mVectorUnit{b}$ exactly cancels against part of the tension, leaving no magnetic force along the field itself. The ratio of thermal to magnetic pressure is the plasma-$\mPlasmaBeta$, and the Alfv\'en Mach number $\mAlfvenMach$ compares this same tension to inertia.

\TODO{Sanity check the sign and prefactor conventions against your preferred reference.}

\end{document}
